\subsection{极大序域上的代数}

设 R 是一个极大序域（VI, §2, No.~5, 第 25 页）。设 C 是型为 \((-1,0)\) 的二次 R-代数；若 \((1,i)\) 是它的典范基，则我们有 \(i^{2}=-1\)。此外，C 是 R 的一个代数闭包（VI, §2, No.~6, 第 26 页，定理 3）。设 H 是型为 \((-1,0,-1)\) 的四元数 R-代数。H 在其典范基 \((1,i,j,k)\) 下的乘法表由下式给出

\[
i ^ {2} = j ^ {2} = k ^ {2} = - 1, \quad i j = - j i = k, \quad - i k = k i = j, \quad j k = - k j = i.
\]

我们把 C 等同于 H 的子代数 \(R + Ri\)。H 的元素 \(q = x + yi + zj + tk\) 的共轭是 \(\overline{q} = x - yi - zj - tk\)。q 的凯莱迹与凯莱范数由下式给出

\[
\mathrm{T} (q) = q + \overline {{q}} = 2 x, \quad \mathrm{N} (q) = q \overline {{q}} = x ^ {2} + y ^ {2} + z ^ {2} + t ^ {2}.
\]

由于 R 是一个序域，我们有 \(N(q) > 0\)（若 \(q \neq 0\)），故 H 是一个体，其中心为 R（VIII, 第 363 页，命题 2）。H 的元素 q 的约化迹与约化范数分别是 \(T(q)\) 与 \(N(q)\)。

定理 1. —— 设 D 是一个有限次数的 R-代数且是一个体。则 D 同构于 R、C 或 H。

把 D 的中心记作 Z，并设 L 是 D 的一个极大交换子域。由 VIII, 第 265 页，推论 2，我们有 \([D:Z]=[L:Z]^{2}\)；由于 C 是 R 的代数闭包，我们还有 \([L:R]\leqslant2\)。因而有三种可能的情形：

\begin{enumerate}
\def\labelenumi{\alph{enumi})}
\item
  我们有 \(\mathrm{R} = \mathrm{Z} = \mathrm{L}\)，故 \([\mathrm{D}:\mathrm{Z}] = 1\) 且 \(\mathrm{D} = \mathrm{R}\)。
\item
  我们有 \(\mathrm{R} \neq \mathrm{Z}\) 且 \(\mathrm{Z} = \mathrm{L}\)，故 \([\mathrm{D} : \mathrm{Z}] = 1\) 且 \(\mathrm{D} = \mathrm{L}\)。在这种情形，\(\mathrm{D}\) 同构于 \(\mathrm{C}\)。
\item
  我们有 R = Z 且 \([L : R] = 2\)，故 \([D : R] = 4\)。由 VIII, 第 364 页，命题 4，R-代数 D 同构于一个型为 \((\alpha, 0, \gamma)\) 的四元数代数，其中 \(\alpha\) 与 \(\gamma\) 是 R 的非零元素。设 \(i \in D - Z\) 满足 \(i^{2} = \alpha\)。我们有 \(\alpha \neq 0\)。若 \(\alpha > 0\)，则存在一个 \(a \in R\) 使得 \(a^{2} = \alpha\)（VI, §2, No.~6, 第 26 页，定理 3）；于是我们有 \((a - i)(a + i) = 0\)，由于 D 是一个体，这是荒谬的。所以我们有 \(\alpha < 0\)。不等式 \(\gamma < 0\) 类似地证明。于是存在 \(R^{*}\) 的元素 a 与 c 使得 \(\alpha = -a^{2}\) 且 \(\gamma = -c^{2}\)（出处同前）。因此代数 D 同构于型为 \((-1,0,-1)\) 的四元数代数（VIII, 第 362 页，注 2），即同构于 H。
\end{enumerate}

注. —— 1) 设 O 是 R 上型为 \((-1,0,-1,-1)\) 的八元数代数（III, 附录， No.~3, 第 615 页）。设 D 是 R 上的一个交错凯莱代数，使得 D 的每个非零元素都有逆元。我们可以证明（VIII, 第 370 页，习题 5）D 同构于 R、C、H 或 O。

*2) 上述结果适用于实数体 \(\mathbf{R}\)。每个是体的有限次数 \(\mathbf{R}\) -代数都同构于 \(\mathbf{R}\)、\(\mathbf{C}\) 或 \(\mathbf{H}\)。

\begin{enumerate}
\def\labelenumi{\arabic{enumi})}
\setcounter{enumi}{2}
\tightlist
\item
  设 A 是体 R 上的一个赋范代数。设 A 是一个体。则 A 同构于 R、C 或 H（``盖尔范德--马祖尔定理''）（参看 Comm. Alg., VI, §6, No.~4, 第 407 页，定理 1 与 TS, I, §2, n° 5, 第 26 页，推论 2）。*
\end{enumerate}

